% \section{Results}

% We present preliminary results from evaluating fifteen classic short stories through our three-layer framework. The human-authored reference texts span 1013 to 8109 words with diverse narrative characteristics. For each story, we generated LLM-rewritten versions using Gemini 2.5 Flash Lite at three temperature settings ($\tau \in \{0.3, 0.8, 1.2\}$), with five independent runs per configuration to account for stochastic variability. The evaluation metrics aggregate results across all fifteen stories, computing means and standard deviations to characterize population-level patterns. These findings represent initial observations that will require further validation through expanded story collections and additional model comparisons.

% \subsection{Surface Linguistic Features}

% Table~\ref{tab:layer1_results} summarizes Layer 1 metrics measuring lexical diversity, syntactic complexity, and stylistic variation. The results reveal consistent patterns across temperature settings. LLM-generated texts exhibit substantially higher lexical diversity compared to human baselines, with MTLD scores ranging from 112.7 at $\tau=0.3$ to 125.0 at $\tau=1.2$, representing increases of 27\% to 41\% over the human baseline of 88.8. This pattern holds across complementary diversity measures including vocd-D and type-token ratios, suggesting LLM outputs employ broader vocabularies regardless of generation temperature.

% \begin{table}[t]
% \centering
% \caption{Layer 1 evaluation results: Surface linguistic features across 15 stories. Values show means; N indicates sample size.}
% \label{tab:layer1_results}
% \small
% \begin{tabular}{lrrrr}
% \hline
% \textbf{Metric} & \textbf{Human} & \textbf{T=0.3} & \textbf{T=0.8} & \textbf{T=1.2} \\
% \hline
% MTLD (Lexical Diversity) & 88.82 & 112.66 & 119.07 & 124.98 \\
% vocd-D & 103.30 & 114.88 & 119.02 & 121.60 \\
% Avg Tree Depth & 4.52 & 4.07 & 4.06 & 4.10 \\
% Avg Clause Depth & 0.25 & 0.17 & 0.17 & 0.17 \\
% Sentence CV & 0.63 & 0.58 & 0.57 & 0.57 \\
% Dialogue Density & 0.004 & 0.005 & 0.005 & 0.005 \\
% \hline
% N (stories) & 15 & 15 & 15 & 15 \\
% \hline
% \end{tabular}
% \end{table}

% Syntactic complexity metrics present an opposing pattern. Human texts achieve greater tree depth (4.52 versus approximately 4.1 for LLM variants) and substantially deeper clause nesting (0.25 versus 0.17), indicating that reference narratives employ more elaborate syntactic structures with greater levels of embedding. The coefficient of variation in sentence length provides a third dimension of contrast, with human texts showing higher variability (0.63) compared to more uniform LLM sentence distributions (0.57--0.58). This suggests human authors deliberately vary sentence rhythm for rhetorical effect, while LLM outputs maintain more consistent pacing. Dialogue density remains comparable across conditions, suggesting this surface feature transfers reliably during controlled rewriting.

% \subsection{Narrative Structure}

% Layer 2 metrics assess coherence, temporal consistency, and plot dynamics at the discourse level. Table~\ref{tab:layer2_results} presents key structural indicators. Entity tracking reveals minimal differences, with both human and LLM texts maintaining approximately 50 named entities and similar persistence patterns around 1.1 mentions per entity on average. This suggests the entity grid framework successfully captures coherence at a coarse level, but may lack sensitivity to subtle differences in how entities participate in narrative progression.

% \begin{table}[t]
% \centering
% \caption{Layer 2 evaluation results: Narrative structure across 15 stories}
% \label{tab:layer2_results}
% \small
% \begin{tabular}{lrrrr}
% \hline
% \textbf{Metric} & \textbf{Human} & \textbf{T=0.3} & \textbf{T=0.8} & \textbf{T=1.2} \\
% \hline
% Entity Count & 49.73 & 49.23 & 49.28 & 49.28 \\
% Entity Persistence & 1.13 & 1.10 & 1.10 & 1.10 \\
% Past Tense (\%) & 0.53 & 0.56 & 0.56 & 0.56 \\
% Sentiment Curve Length & 207.13 & 147.16 & 153.27 & 153.92 \\
% Climax Position & 0.45 & 0.49 & 0.52 & 0.51 \\
% \hline
% N (stories) & 15 & 15 & 15 & 15 \\
% \hline
% \end{tabular}
% \end{table}

% Temporal markers show slight increases in past tense usage in LLM texts (56\% versus 53\% in human references), potentially reflecting differences in narrative voice or tendency toward more straightforward chronological progression. More substantial divergence appears in plot tension dynamics. Sentiment curve length, measuring the cumulative magnitude of emotional trajectory fluctuations, averages 207.1 in human stories but only 147.2 to 153.9 in LLM variants. This 29\% reduction suggests LLM narratives exhibit less emotional volatility, producing smoother affective arcs with fewer sharp transitions between positive and negative sentiment states. Climax position shifts slightly later in LLM texts (0.49--0.52 versus 0.45), indicating a tendency toward more delayed narrative peaks compared to the human preference for earlier dramatic culmination.

% \subsection{Deep Semantic Properties}

% Layer 3 employs embedding-based methods to assess thematic structure and emotional granularity. Table~\ref{tab:layer3_results} shows results for topic modeling and semantic consistency metrics. Topic diversity remains remarkably stable at 0.045 across all conditions, suggesting that high-level thematic structure transfers effectively during rewriting. The number of distinct topics identified by BERTopic clustering does not vary meaningfully between human and LLM texts, indicating preservation of broad semantic organization.

% \begin{table}[t]
% \centering
% \caption{Layer 3 evaluation results: Deep semantic properties across 15 stories}
% \label{tab:layer3_results}
% \small
% \begin{tabular}{lrrrr}
% \hline
% \textbf{Metric} & \textbf{Human} & \textbf{T=0.3} & \textbf{T=0.8} & \textbf{T=1.2} \\
% \hline
% Topic Diversity & 0.045 & 0.045 & 0.045 & 0.045 \\
% Thematic Consistency & 0.46 & 0.50 & 0.49 & 0.50 \\
% Semantic Drift & 0.47 & 0.46 & 0.45 & 0.42 \\
% Emotion Diversity & 5.47 & 4.97 & 5.09 & 4.99 \\
% \hline
% N (stories) & 15 & 15 & 15 & 15 \\
% \hline
% \end{tabular}
% \end{table}

% Thematic consistency, measured through cosine similarity between consecutive narrative segments in embedding space, shows slightly higher values for LLM texts (0.49--0.50 versus 0.46 for humans). This increased consistency could indicate either superior coherence or reduced thematic development, depending on whether variation serves narrative progression. Semantic drift, measuring similarity decay over narrative distance, decreases at higher temperatures (0.42 at $\tau=1.2$ versus 0.47 for humans), potentially reflecting more stable semantic trajectories in LLM outputs. Emotional granularity assessed through 27-category emotion classification reveals modest reductions in LLM texts, with 5.0 distinct emotions on average compared to 5.5 in human references. This compression of emotional range aligns with the reduced sentiment volatility observed in Layer 2, suggesting LLM narratives operate within a narrower affective space.

% \subsection{Preliminary Interpretation}

% These results establish that current LLM-generated narratives achieve surface-level linguistic elaboration through expanded vocabularies while exhibiting reduced complexity in syntactic structure, emotional dynamics, and stylistic variation. The pattern suggests a tradeoff where lexical richness comes at the expense of structural depth and rhetorical sophistication. Human texts employ deliberate syntactic complexity and sentence rhythm variation as expressive devices, features that LLM outputs achieve less consistently despite their broader lexical range. The flattened emotional trajectories and reduced affective granularity in LLM narratives point to challenges in modeling the dramatic pacing and psychological nuance characteristic of literary fiction.

% Several caveats apply to these preliminary findings. The evaluation corpus comprises only fifteen stories from a narrow historical period, limiting generalization to contemporary or diverse narrative styles. The single-model comparison (Gemini 2.5 Flash Lite) does not capture potential variations across architectural families or training regimes. The paired evaluation design, while controlling for plot structure, may inadvertently favor metrics insensitive to higher-order narrative qualities. Future work will expand the story collection, incorporate multiple LLM families, and develop additional metrics targeting narrative elements our current framework may underweight, such as character development, thematic subtlety, and stylistic distinctiveness.
